%==============================================================================
%  RELATÓRIO TÉCNICO — BALL BALANCER 3RPS (ESP32-S3)
%  Compilar com: pdflatex (2x para sumário/referências)  ou  Overleaf
%==============================================================================
\documentclass[11pt,a4paper,oneside]{article}

%------------------------------- Codificação ---------------------------------
\usepackage[utf8]{inputenc}
\usepackage[T1]{fontenc}
\usepackage[brazilian]{babel}
\usepackage{lmodern}
\usepackage{microtype}

%------------------------------- Layout --------------------------------------
\usepackage[a4paper,left=2.6cm,right=2.6cm,top=2.8cm,bottom=2.6cm]{geometry}
\usepackage{setspace}
\onehalfspacing
\usepackage{parskip}

%------------------------------- Gráficos ------------------------------------
\usepackage{graphicx}
\usepackage{xcolor}
\usepackage{tikz}
\usetikzlibrary{shapes.geometric,arrows.meta,positioning,calc,backgrounds}

%------------------------------- Matemática ----------------------------------
\usepackage{amsmath}
\usepackage{amssymb}
\usepackage{mathtools}

%------------------------------- Tabelas -------------------------------------
\usepackage{booktabs}
\usepackage{array}
\usepackage{colortbl}
\usepackage{multirow}
\usepackage{longtable}
\usepackage{tabularx}
\usepackage{ragged2e}

%------------------------------- Estrutura -----------------------------------
\usepackage{enumitem}
\setlist{itemsep=2pt,topsep=3pt}
\usepackage{caption}
\captionsetup{font=small,labelfont=bf,labelsep=endash}
\usepackage[most]{tcolorbox}
\usepackage{fancyhdr}
\usepackage{titlesec}
\usepackage{listings}

% fontawesome5 é opcional: se não existir no compilador, os ícones viram vazio
% e o documento compila do mesmo jeito.
\IfFileExists{fontawesome5.sty}{%
  \usepackage{fontawesome5}%
}{%
  \providecommand{\faExclamationTriangle}{}%
  \providecommand{\faCheckCircle}{}%
}

%------------------------------- Hyperlinks ----------------------------------
\usepackage[colorlinks=true]{hyperref}

%==============================================================================
%  PALETA DE CORES
%==============================================================================
\definecolor{primary}{HTML}{1A4E8A}     % azul profundo
\definecolor{secondary}{HTML}{0E7C7B}   % verde-petróleo
\definecolor{accent}{HTML}{E8833A}      % laranja
\definecolor{danger}{HTML}{C0392B}      % vermelho
\definecolor{success}{HTML}{2E8B57}     % verde
\definecolor{inkdark}{HTML}{1C2A39}     % quase-preto
\definecolor{paperbg}{HTML}{F7F9FB}     % fundo claro
\definecolor{rowgray}{HTML}{EEF3F7}     % linha de tabela
\definecolor{codebg}{HTML}{F4F6F8}

\hypersetup{
  linkcolor=primary, citecolor=secondary, urlcolor=accent,
  pdftitle={Relatório Técnico — Ball Balancer 3RPS},
  pdfauthor={Engenharia Elétrica — UNIFOR}
}

%==============================================================================
%  CABEÇALHO E RODAPÉ
%==============================================================================
\pagestyle{fancy}
\fancyhf{}
\renewcommand{\headrulewidth}{0.4pt}
\renewcommand{\footrulewidth}{0.4pt}
\fancyhead[L]{\small\color{primary}\textbf{Ball Balancer 3RPS}}
\fancyhead[R]{\small\color{inkdark}ESP32-S3 · Controle PID}
\fancyfoot[L]{\small\color{inkdark}Engenharia Elétrica — UNIFOR}
\fancyfoot[R]{\small\thepage}

%==============================================================================
%  ESTILO DE SEÇÕES
%==============================================================================
\titleformat{\section}
  {\Large\bfseries\color{primary}}
  {\color{accent}\thesection}{0.7em}{}
  [{\color{primary!30}\titlerule[1pt]}]
\titlespacing*{\section}{0pt}{1.4em}{0.7em}
\titleformat{\subsection}
  {\large\bfseries\color{secondary}}
  {\color{secondary}\thesubsection}{0.6em}{}
\titleformat{\subsubsection}
  {\normalsize\bfseries\color{inkdark}}
  {\color{inkdark!60}\thesubsubsection}{0.6em}{}

%==============================================================================
%  CAIXAS PERSONALIZADAS (tcolorbox)
%==============================================================================
% Caixa de destaque/informação
\newtcolorbox{infobox}[1][]{
  enhanced, colback=paperbg, colframe=primary, boxrule=0.8pt,
  arc=2pt, left=8pt, right=8pt, top=6pt, bottom=6pt,
  fonttitle=\bfseries\color{white}, coltitle=white,
  attach boxed title to top left={xshift=8pt,yshift=-3pt},
  boxed title style={colback=primary,arc=1pt}, title={#1}
}
% Caixa de problema
\newtcolorbox{problembox}[1][Problema]{
  enhanced, colback=danger!4, colframe=danger, boxrule=0.8pt,
  arc=2pt, left=8pt, right=8pt, top=6pt, bottom=6pt,
  fonttitle=\bfseries\color{white}, coltitle=white,
  attach boxed title to top left={xshift=8pt,yshift=-3pt},
  boxed title style={colback=danger,arc=1pt},
  title={\faExclamationTriangle\;\,#1}
}
% Caixa de solução
\newtcolorbox{solutionbox}[1][Solução adotada]{
  enhanced, colback=success!4, colframe=success, boxrule=0.8pt,
  arc=2pt, left=8pt, right=8pt, top=6pt, bottom=6pt,
  fonttitle=\bfseries\color{white}, coltitle=white,
  attach boxed title to top left={xshift=8pt,yshift=-3pt},
  boxed title style={colback=success,arc=1pt},
  title={\faCheckCircle\;\,#1}
}
% Caixa de equação-chave
\newtcolorbox{keybox}{
  enhanced, colback=secondary!6, colframe=secondary, boxrule=1pt,
  arc=3pt, left=10pt, right=10pt, top=8pt, bottom=8pt, halign=center
}

%==============================================================================
%  ESTILO DE CÓDIGO
%==============================================================================
\lstdefinestyle{custom}{
  basicstyle=\ttfamily\footnotesize,
  backgroundcolor=\color{codebg},
  keywordstyle=\color{primary}\bfseries,
  commentstyle=\color{secondary}\itshape,
  stringstyle=\color{accent},
  numberstyle=\tiny\color{gray},
  frame=leftline, framerule=2pt, rulecolor=\color{primary},
  breaklines=true, showstringspaces=false,
  xleftmargin=10pt, aboveskip=8pt, belowskip=8pt,
  numbers=none
}
\lstset{style=custom}

% Comando para parâmetros inline — caixinha arredondada e discreta
\newtcbox{\param}{on line, boxrule=0.4pt, boxsep=0pt,
  left=3pt, right=3pt, top=1.5pt, bottom=1.5pt, arc=2pt,
  colback=rowgray, colframe=primary!25,
  fontupper=\ttfamily\small\color{inkdark}}
% Linha de tabela zebrada
\newcommand{\rowstyle}{\rowcolor{rowgray}}

%==============================================================================
\begin{document}

%==============================================================================
%  CAPA
%==============================================================================
\begin{titlepage}
\thispagestyle{empty}
\begin{tikzpicture}[remember picture,overlay]
  % --- Banner superior sólido ---
  \fill[primary] (current page.north west) rectangle ([yshift=-6cm]current page.north east);
  % faixa de acento fina abaixo do banner
  \fill[accent]  ([yshift=-6cm]current page.north west) rectangle ([yshift=-6.18cm]current page.north east);
  % --- Rodapé sólido ---
  \fill[primary] (current page.south west) rectangle ([yshift=2.2cm]current page.south east);
  \fill[accent]  ([yshift=2.2cm]current page.south west) rectangle ([yshift=2.32cm]current page.south east);

  % --- Texto do banner ---
  \node[anchor=north west, text=white, align=left]
    at ([xshift=2.6cm,yshift=-1.5cm]current page.north west)
    {{\fontsize{15}{18}\selectfont\textsc{Universidade de Fortaleza — UNIFOR}}\\[3pt]
     {\fontsize{11}{14}\selectfont Engenharia Elétrica \;\textbar\; Controle Digital}};
  \node[anchor=south east, text=white]
    at ([xshift=-2.6cm,yshift=-5.6cm]current page.north east)
    {\fontsize{10}{12}\selectfont\itshape Relatório Técnico};

  % --- Texto do rodapé ---
  \node[anchor=west, text=white]
    at ([xshift=2.6cm,yshift=1.1cm]current page.south west)
    {\small Fortaleza · \today};
  \node[anchor=east, text=white]
    at ([xshift=-2.6cm,yshift=1.1cm]current page.south east)
    {\small Grupo de 8 integrantes};
\end{tikzpicture}

% --- Bloco central do título ---
\vspace*{7.5cm}
\begin{center}
  {\fontsize{13}{15}\selectfont\bfseries\color{accent}\textsc{Relatório Técnico}}\\[14pt]
  {\fontsize{46}{50}\selectfont\bfseries\color{primary}Ball Balancer 3RPS}\\[14pt]
  {\color{secondary}\rule{6cm}{2.5pt}}\\[14pt]
  {\Large\color{inkdark}Mesa Equilibrista com Controle PID Duplo}\\[8pt]
  {\large\color{inkdark!70}Cinemática Inversa 3RPS \;·\; ESP32-S3 \;·\; Três Motores NEMA 17}
\end{center}

\vfill

\begin{center}
\begin{tcolorbox}[enhanced,width=0.86\textwidth,colback=paperbg,
  colframe=primary,boxrule=1pt,arc=5pt,
  left=14pt,right=14pt,top=10pt,bottom=10pt]
  {\bfseries\color{primary}Objeto do relatório}\par\smallskip
  {\small Documentar de forma reproduzível o projeto de uma mesa equilibrista
  controlada por PID digital — fundamentos de controle, arquitetura de
  hardware e firmware, dificuldades enfrentadas e levantamento de custos —
  servindo de referência para futuras práticas.}
\end{tcolorbox}
\end{center}

\vspace{2.4cm}
\end{titlepage}

%==============================================================================
%  RESUMO
%==============================================================================
\section*{Resumo}
\addcontentsline{toc}{section}{Resumo}

Este relatório documenta o projeto de uma \textbf{mesa equilibrista} (\emph{ball
balancer}): uma plataforma cuja inclinação é controlada para manter uma esfera
no centro, apesar de perturbações. O sistema usa um microcontrolador
\textbf{ESP32-S3}, um \textbf{sensor de posição por tela resistiva de 4 fios} e
\textbf{três motores de passo NEMA 17} dispostos em arranjo cinemático
\textbf{3RPS} (\emph{Revolute–Prismatic–Spherical}).

O controle é feito por \textbf{dois laços PID independentes} (um para cada eixo,
X e Y), executados a \textbf{100\,Hz} no firmware, com derivativo filtrado e
anti-\emph{windup}. A saída de cada PID é uma inclinação desejada da mesa, que a
\textbf{cinemática inversa 3RPS} traduz em três ângulos de motor; cada motor é
acionado por um perfil de velocidade trapezoidal. A fundamentação teórica
(modelo da planta, função de transferência, projeto dos ganhos por alocação de
polos e discretização) baseia-se no TCC de Cordeiro (UNESP, 2022), adaptando a
solução de \textbf{dois servos} do autor para os \textbf{três motores} deste rig.

São relatados os problemas práticos enfrentados — pinagem da tela, fabricação da
PCB, fragilidade da estrutura impressa em 3D, acesso a componentes mecânicos — e
suas soluções, além de um levantamento de custos real do grupo.

\medskip
\noindent\textbf{Palavras-chave:} Controle PID · Duplo Integrador · Cinemática
3RPS · Motor de Passo · Tela Resistiva · Sistema Embarcado · ESP32.

%==============================================================================
%  SUMÁRIO
%==============================================================================
\newpage
{\hypersetup{linkcolor=inkdark}
\renewcommand{\contentsname}{\color{primary}Sumário}
\tableofcontents}

\newpage

%==============================================================================
%  1. INTRODUÇÃO
%==============================================================================
\section{Introdução}

\subsection{Contexto}

A mesa equilibrista é um problema clássico de controle e um excelente banco de
ensino: reúne, num único equipamento, uma planta instável, restrições reais de
atuação e a necessidade de realimentação rápida. Equilibrar uma esfera sobre uma
superfície inclinável exige um controlador que reaja à posição \emph{e} à
velocidade da bola — o que torna a demonstração do papel de cada termo (P, I, D)
particularmente didática.

\begin{infobox}[Por que este projeto é instrutivo]
\begin{itemize}
  \item \textbf{Planta naturalmente instável:} a bola num plano é um
  \emph{duplo integrador}, marginalmente estável — sem o termo derivativo, o
  laço oscila para sempre.
  \item \textbf{Restrições físicas reais:} saturação da inclinação, cinemática
  não-linear, ruído de sensor e perda de passos nos motores.
  \item \textbf{Dois eixos:} dois PID idênticos e independentes (X e Y),
  acoplados apenas pela geometria 3RPS.
  \item \textbf{Ajuste ao vivo:} ganhos e calibração persistem em memória
  não-volátil (NVS) e são ajustados sem recompilar.
\end{itemize}
\end{infobox}

\subsection{Objetivos}

\begin{enumerate}
  \item Implementar um controlador PID digital duplo embarcado, executando em
  malha fechada a 100\,Hz.
  \item Resolver a cinemática inversa 3RPS para converter a inclinação desejada
  em três ângulos de motor.
  \item Integrar o sensor resistivo com calibração e orientação ajustáveis em
  tempo de execução.
  \item Validar o controlador por simulação (sem hardware) e, em seguida, no
  rig real.
  \item Produzir documentação reproduzível para futuras turmas.
\end{enumerate}

\subsection{Organização do documento}

\begin{center}
\small
\begin{tabularx}{\textwidth}{>{\bfseries\color{primary}}l X}
\toprule
Seção 2 & Referências de projeto (Aaed Musa e TCC de Cordeiro). \\
\rowstyle Seção 3 & Fundamentação teórica de controle (modelo, PID, discretização). \\
Seção 4 & Arquitetura de hardware e pinagem real. \\
\rowstyle Seção 5 & Firmware, ferramentas e protocolo serial. \\
Seção 6 & Problemas enfrentados e soluções. \\
\rowstyle Seção 7 & Materiais e custos reais do grupo. \\
Seção 8 & Validação e resultados. \\
\rowstyle Seção 9 & Pendências e trabalhos futuros. \\
\bottomrule
\end{tabularx}
\end{center}

%==============================================================================
%  2. REFERÊNCIAS DE PROJETO
%==============================================================================
\section{Referências de Projeto}

\subsection{Ball Balancer de Aaed Musa (arquitetura mecânica)}

A arquitetura mecânica e a cinemática inversa 3RPS deste projeto baseiam-se no
\emph{Ball Balancer} de Aaed Musa (Instructables). Dela vêm:

\begin{itemize}
  \item o arranjo \textbf{3RPS} com três motores defasados de 120°;
  \item a função fechada de cinemática inversa (\texttt{machine\_theta}),
  portada fielmente para C em \texttt{main/kinematics.c};
  \item o uso de motores de passo NEMA 17 com drivers STEP/DIR.
\end{itemize}

\subsection{TCC de Cordeiro — UNESP, 2022 (base teórica de controle)}

A teoria de controle apoia-se no trabalho de graduação de \textbf{Giovanna
Ferreira Cordeiro} (\emph{Controle Clássico da Planta Ball Balancer}, UNESP/FEIS,
2022), que projeta um controlador em \textbf{cascata} para uma mesa equilibrista
do tipo \textbf{Quanser 2 DOF (dois servomotores)}:

\begin{itemize}
  \item \textbf{Malha externa} ($C_{bb}$, PID de posição): do erro de posição
  calcula a inclinação desejada da mesa;
  \item \textbf{Malha interna} ($C_s$, PD rápido): leva cada servo ao ângulo
  desejado;
  \item \textbf{Projeto de ganhos:} alocação de polos para
  \emph{overshoot}~$\le 7{,}5\%$ e tempo de acomodação~$\le 2{,}5$\,s.
\end{itemize}

\begin{infobox}[Diferença essencial: 2 servos $\rightarrow$ 3 motores]
O TCC usa \textbf{dois servos} (um ângulo por eixo). Nosso rig usa \textbf{três
motores de passo} em 3RPS. A adaptação preserva a \emph{malha externa} de
posição (mesma teoria, mesmos ganhos) e troca a \emph{malha interna}: em vez de
um servo por eixo, a inclinação desejada $(n_x, n_y)$ é distribuída entre
\textbf{três ângulos} $(\theta_A,\theta_B,\theta_C)$ pela cinemática inversa
3RPS. Como essa cinemática é um \emph{mapa estático não-linear}, a análise
linear de controle permanece válida em primeira ordem.
\end{infobox}

%==============================================================================
%  3. FUNDAMENTAÇÃO TEÓRICA
%==============================================================================
\section{Fundamentação Teórica de Controle}

\subsection{Modelo da planta: a bola num plano inclinado}

Para uma esfera maciça de raio $r$ e massa $m$ rolando \emph{sem deslizar} num
plano inclinado de ângulo $\alpha$, somam-se a 2ª lei de Newton (translação) e a
equação de Euler (rotação):

\begin{equation}
  m\,\ddot{x} = m g \sin\alpha - F_{at}, \qquad
  I\,\dot\omega = F_{at}\,r, \qquad \dot{x} = \omega r .
\end{equation}

Com $I = \tfrac{2}{5}m r^2$ (esfera maciça), elimina-se a força de atrito:

\begin{equation}
  \left(m + \frac{I}{r^2}\right)\ddot{x} = m g \sin\alpha
  \;\Longrightarrow\; \tfrac{7}{5}\,\ddot{x} = g\sin\alpha .
\end{equation}

\begin{keybox}
$\displaystyle \ddot{x} = \frac{5}{7}\, g \sin\alpha \;\approx\; \frac{5}{7}\, g\,\alpha$
\end{keybox}

A fração $5/7\approx 0{,}714$ é a constante \param{ROLL\_FACTOR} da simulação. A
aproximação $\sin\alpha\approx\alpha$ vale na faixa de operação (a inclinação é
limitada a $|u|\le 0{,}25$\,rad $\approx 14{,}3°$).

\subsubsection{Função de transferência}

A saída do PID, $u=n_x$, é a componente $x$ do vetor normal unitário da mesa;
para pequenos ângulos $n_x\approx\tan\alpha\approx\alpha$. Com $g=9810$\,mm/s²:

\begin{equation}
  \ddot{x} = \underbrace{\tfrac{5}{7}\,g}_{A}\,u,
  \qquad A = \tfrac{5}{7}\times 9810 \approx 7007\ \text{mm/s}^2 .
\end{equation}

Aplicando Laplace (condições iniciais nulas):

\begin{keybox}
$\displaystyle P(s) = \frac{X(s)}{U(s)} = \frac{A}{s^2}, \qquad A \approx 7007\ \text{mm/s}^2$
\end{keybox}

A planta é um \textbf{duplo integrador} (tipo 2). Consequências de projeto:

\begin{center}
\small
\begin{tabularx}{\textwidth}{>{\bfseries}l X}
\toprule
Erro estacionário & Nulo a degrau e a rampa de referência. \\
\rowstyle Polos & Os dois ficam na origem ($s=0$): \emph{marginalmente estável}. \\
Implicação & Realimentação só proporcional $\Rightarrow$ oscilação não-amortecida.
O termo \textbf{derivativo é obrigatório}. \\
\bottomrule
\end{tabularx}
\end{center}

\subsection{O controlador PID}

\begin{equation}
  u(t) = K_p\,e(t) + K_i\!\int_0^t\! e(\tau)\,d\tau + K_d\,\dot{e}(t),
  \qquad e = x - r .
\end{equation}

No firmware (\texttt{main/pid.c}) o derivativo é \textbf{filtrado} por um
passa-baixa de 1ª ordem ($\tau_d = 0{,}03$\,s) para não amplificar o ruído do
sensor:

\begin{equation}
  C(s) = K_p + \frac{K_i}{s} + \frac{K_d\,s}{1 + \tau_d\,s}.
\end{equation}

\begin{center}
\small
\renewcommand{\arraystretch}{1.2}
\begin{tabularx}{\textwidth}{>{\bfseries\color{primary}}c >{\RaggedRight}X >{\RaggedRight}X}
\toprule
Termo & \normalfont\bfseries Função nesta planta & \normalfont\bfseries Risco se exagerado \\
\midrule
P \,($K_p$) & Define a rigidez/velocidade da resposta ($\omega_n$). Sozinho, oscila. & Oscilação crescente. \\
\rowstyle D \,($K_d$) & \textbf{Essencial:} adiciona amortecimento (freia pela velocidade da bola). & Lentidão; sensível a ruído. \\
I \,($K_i$) & Remove desvio residual (mesa fora de nível / assimetria mecânica). & \emph{Windup}; instabilidade lenta. \\
\bottomrule
\end{tabularx}
\end{center}

\subsection{Malha fechada e analogia de 2ª ordem}

Desprezando o filtro do D e o termo I (pequeno), a malha fechada PD sobre
$A/s^2$ tem equação característica:

\begin{equation}
  s^2 + A K_d\,s + A K_p = 0 .
\end{equation}

Comparando com a forma canônica $s^2 + 2\zeta\omega_n s + \omega_n^2$:

\begin{keybox}
$\displaystyle \omega_n = \sqrt{A\,K_p} \qquad\qquad
   \zeta = \frac{K_d}{2}\sqrt{\frac{A}{K_p}}$
\end{keybox}

onde $\omega_n$ indica \emph{quão rápido} a bola converge e $\zeta$ governa o
\emph{overshoot} ($\zeta<1$ oscila, $\zeta=1$ crítico, $\zeta>1$ sobreamortecido);
o tempo de acomodação a 2\% é $t_s\approx 4/(\zeta\omega_n)$.

\subsection{Projeto dos ganhos por alocação de polos}

Seguindo o TCC, fixam-se os requisitos \emph{overshoot}~$\le 7{,}5\%$ e
$t_s\le 2{,}5$\,s. Pelas fórmulas de 2ª ordem (Ogata):

\begin{equation}
  \zeta = \sqrt{\frac{\ln^2(PO/100)}{\pi^2 + \ln^2(PO/100)}} \approx 0{,}636,
  \qquad
  \omega_n \approx 2{,}19\ \text{rad/s}.
\end{equation}

Como a malha fechada é de 3ª ordem (com o I), alocam-se os polos como um par
dominante mais um polo real $p_o$. Igualando os coeficientes de
$s^3 + A K_d s^2 + A K_p s + A K_i = 0$ ao polinômio desejado (com $p_o = 1$):

\begin{equation}
  A K_d = 2\zeta\omega_n + p_o, \qquad
  A K_p = \omega_n^2 + 2\zeta\omega_n p_o, \qquad
  A K_i = \omega_n^2 p_o .
\end{equation}

\begin{center}
\small
\begin{tabularx}{0.9\textwidth}{>{\bfseries\color{primary}}c >{\centering}X >{\centering\arraybackslash}X}
\toprule
Ganho & \normalfont\bfseries Fórmula & \normalfont\bfseries Valor (PID inicial) \\
\midrule
$K_p$ & $(\omega_n^2 + 2\zeta\omega_n p_o)/A$ & $1{,}08\times10^{-3}$ \\
\rowstyle $K_i$ & $\omega_n^2 p_o/A$ & $6{,}83\times10^{-4}$ \\
$K_d$ & $(2\zeta\omega_n + p_o)/A$ & $5{,}40\times10^{-4}$ \\
\bottomrule
\end{tabularx}
\end{center}

\subsubsection{``PID melhorado'' (amortecido)}

O PID inicial dá \emph{overshoot} de $25$–$30\%$. Como no TCC, reduz-se $K_i$ e
aumenta-se $K_d$ (alvo $\zeta\approx 1{,}15$). Valores adotados na simulação
\texttt{PIDSimba.py}:

\begin{keybox}
$K_p = 1{,}08\times10^{-3} \qquad K_i = 3{,}0\times10^{-4} \qquad K_d = 9{,}0\times10^{-4}$
\end{keybox}

Com eles a bola converge dos quatro cantos, \emph{overshoot} $\approx 15\%$ e
erro residual final $\approx 1$\,mm (tendendo a zero pela ação integral).

\subsubsection{Ganhos padrão do firmware — conservadores e por quê}

O firmware (\texttt{main/control.c}) parte de valores propositalmente
sobreamortecidos:

\begin{keybox}
$K_p = 8{,}0\times10^{-4} \qquad K_i = 2{,}0\times10^{-5} \qquad K_d = 1{,}2\times10^{-2}$
\end{keybox}

\begin{enumerate}
  \item \textbf{Saturação proposital do D:} com a saída limitada a $\pm0{,}25$, o
  termo D sozinho satura a partir de $v\approx 0{,}25/0{,}012\approx 21$\,mm/s,
  atuando como um \textbf{freio de velocidade quase \emph{bang-bang}} — muito
  eficaz contra empurrões.
  \item \textbf{Ponto de partida seguro:} o ajuste fino é feito \emph{no rig}; os
  defaults precisam ser mansos para não derrubar a bola na primeira energização.
\end{enumerate}

\subsection{Controle digital (discretização)}

O laço roda a $100$\,Hz $\Rightarrow T = 0{,}01$\,s. Nyquist (50\,Hz) é muito
superior a $\omega_n/2\pi\approx 0{,}35$\,Hz, então a amostragem é folgada. (O laço
subiu de 50 para 100\,Hz para reduzir o atraso de amostragem e melhorar a resposta.)

\begin{center}
\small
\renewcommand{\arraystretch}{1.4}
\begin{tabularx}{\textwidth}{>{\bfseries}l >{\centering\arraybackslash}p{6.2cm} >{\RaggedRight\arraybackslash}X}
\toprule
Bloco & \normalfont\bfseries Discretização & \normalfont\bfseries No código \\
\midrule
Integral & $I[k] = I[k-1] + e[k]\,T$ & \texttt{p->integ += err*dt;} \\
\rowstyle Derivada & $d[k] = \dfrac{e[k]-e[k-1]}{T}$ & diferença regressiva \\
Filtro do D & $\alpha = \dfrac{T}{\tau_d+T} = 0{,}25$ & IIR de 1ª ordem \\
\bottomrule
\end{tabularx}
\end{center}

\subsubsection{Anti-\emph{windup} por \emph{back-calculation}}

Ao saturar em $\pm0{,}25$, o integrador é puxado de volta pela quantidade exata
que excedeu, evitando \emph{overshoot} gigante ao sair da saturação:

\begin{equation}
  \text{se } u>u_{max}:\quad I \mathrel{-}= \frac{u-u_{max}}{K_i},\quad u=u_{max}.
\end{equation}

\begin{lstlisting}[language=C,caption={Trecho real de \texttt{main/pid.c} — anti-windup.}]
if (out > p->out_max) {
    if (p->ki != 0.0f) p->integ -= (out - p->out_max) / p->ki;
    out = p->out_max;
}
\end{lstlisting}

A cada borda de ``sem toque'' o controlador é \textbf{resetado}
(\texttt{pid\_reset}), evitando arrastar estado integral/derivativo entre
sessões.

\subsection{Cinemática inversa 3RPS}

O PID produz o vetor de inclinação $(n_x,n_y)$. A função
\texttt{machine\_theta(perna, hz, nx, ny)} resolve, em forma fechada, o ângulo de
cada motor:

\begin{equation}
  (n_x, n_y) \;\xrightarrow{\ \text{cinemática inversa}\ }\; (\theta_A, \theta_B, \theta_C).
\end{equation}

A conversão para passos usa a posição neutra $\theta_{orig}=\texttt{machine\_theta}(\cdot,n_x{=}0,n_y{=}0)$:

\begin{equation}
  \text{passos} = (\theta_{orig} - \theta)\cdot\frac{200\cdot\text{microsteps}}{360°}.
\end{equation}

\subsection{Atuador: motor de passo com perfil trapezoidal}

Cada motor segue o alvo com um \textbf{perfil trapezoidal de velocidade}
(\texttt{main/steppers.c}, timer de 10\,kHz): acelera até a velocidade máxima e
desacelera a tempo de parar \emph{exatamente} no alvo. A velocidade da qual ainda
se freia a tempo é

\begin{equation}
  v_{lim} = \sqrt{2\,a\,d},
\end{equation}

com $d$ = distância restante (passos) e $a$ = aceleração. Sem essa rampa, a
partida brusca causaria perda de passos. Em equilíbrio:
\param{SPEED\_BALANCE = 1800 passos/s} e \param{ACCEL\_BALANCE = 16000 passos/s²}
(suaves; pressupõem o DRV8825 em 1/16).

\subsection{Estabilidade}

\begin{itemize}
  \item \textbf{Planta $A/s^2$ isolada:} marginalmente estável (2 polos na origem).
  \item \textbf{Com PD:} todos os coeficientes da característica são positivos
  $\Rightarrow$ malha estável para quaisquer $K_p,K_d>0$ (Routh–Hurwitz). O D é
  o que garante a estabilidade.
  \item \textbf{Com o I (pequeno):} adiciona um polo lento ($T_i$ longo, 25–40\,s),
  sem reduzir relevantemente a margem de fase.
  \item \textbf{Amostragem:} atraso de $\approx T/2 = 0{,}01$\,s, desprezível
  frente a $1/\omega_n\approx 0{,}7$\,s.
\end{itemize}

%==============================================================================
%  3.X DIAGRAMA DE BLOCOS
%==============================================================================
\subsection{Diagrama de blocos do sistema}

A Figura~\ref{fig:blocos} resume o laço completo (por eixo): o erro alimenta o
PID, cuja saída de inclinação passa pela cinemática 3RPS, é realizada pelos três
motores e move a bola; a tela resistiva fecha a realimentação.

\begin{figure}[h]
\centering
\begin{tikzpicture}[
  font=\small,
  block/.style={draw=primary,fill=primary!8,thick,rounded corners,
    minimum height=1.1cm,minimum width=2.0cm,align=center},
  plant/.style={draw=accent,fill=accent!10,thick,rounded corners,
    minimum height=1.1cm,minimum width=2.0cm,align=center},
  sum/.style={draw=inkdark,fill=white,thick,circle,minimum size=0.7cm,inner sep=0pt},
  arrow/.style={-{Stealth[length=2.5mm]},thick,primary},
  fb/.style={-{Stealth[length=2.5mm]},thick,danger}
]
  \node[sum] (sum) {$+$};
  \node[left=0.7cm of sum] (ref) {$r$ \scriptsize(centro)};
  \node[block,right=0.8cm of sum] (pid) {PID\\\scriptsize$K_p,K_i,K_d$};
  \node[block,right=0.9cm of pid] (kin) {Cinemática\\\scriptsize inversa 3RPS};
  \node[block,right=0.9cm of kin] (mot) {3$\times$ NEMA 17\\\scriptsize trapezoidal};
  \node[plant,right=0.9cm of mot] (plant) {Planta\\\scriptsize $A/s^2$};
  \node[right=0.8cm of plant] (out) {$x$};

  \draw[arrow] (ref) -- (sum);
  \draw[arrow] (sum) -- node[above,scriptsize]{$e$} (pid);
  \draw[arrow] (pid) -- node[above,scriptsize]{$n_x$} (kin);
  \draw[arrow] (kin) -- node[above,scriptsize]{$\theta_{A,B,C}$} (mot);
  \draw[arrow] (mot) -- node[above,scriptsize]{$\alpha$} (plant);
  \draw[arrow] (plant) -- (out);

  % realimentação
  \coordinate (tap) at ($(plant.east)+(0.4,0)$);
  \draw[fb] (tap) |- ($(sum.south)+(0,-1.1)$) -| node[pos=0.25,below,scriptsize]{sensor (tela resistiva)} (sum.south);
  \node[below=1.1cm of sum,scriptsize] {$-$};
\end{tikzpicture}
\caption{Diagrama de blocos do laço de controle (idêntico para X e Y).}
\label{fig:blocos}
\end{figure}

%==============================================================================
%  4. HARDWARE
%==============================================================================
\section{Arquitetura de Hardware}

\subsection{Componentes}

\begin{center}
\small
\renewcommand{\arraystretch}{1.25}
\begin{tabularx}{\textwidth}{>{\bfseries\color{primary}}l >{\RaggedRight\arraybackslash}X}
\toprule
Microcontrolador & ESP32-S3 DevKitC-1 (dual-core, ADC 12 bits). \\
\rowstyle Sensor de posição & Tela resistiva 4 fios — VSDISPLAY VS084TP-A1 (8,4''). \\
Área ativa & $187\times141$\,mm (largura $\times$ altura). \\
\rowstyle Atuadores & 3$\times$ NEMA 17 (1,8°/passo $=$ 200 passos/rev). \\
Drivers de passo & A4988 / DRV8825 / TMC2208 (interface STEP/DIR idêntica). \\
\rowstyle Microstepping & 16$\times$ (\texttt{MICROSTEPS}, via DIP switches). \\
\bottomrule
\end{tabularx}
\end{center}

\subsection{Pinagem (verificada no firmware)}

\subsubsection{Tela resistiva — conector FPC}

Cabo flat com o conector voltado para a \textbf{direita} (olhando a superfície).

\begin{center}
\small
\renewcommand{\arraystretch}{1.2}
\begin{tabular}{>{\bfseries}c c c c}
\toprule
\rowcolor{primary}
\textcolor{white}{Pino FPC} & \textcolor{white}{Sinal} & \textcolor{white}{GPIO} & \textcolor{white}{Papel} \\
\midrule
1 & Y+ & 4 & ADC1\_CH3 — \emph{sense} de X \\
\rowstyle 2 & X$-$ & 5 & ADC1\_CH4 — \emph{sense} de Y \\
3 & Y$-$ & 8 & ADC1\_CH7 — acionado na medida de Y \\
\rowstyle 4 & X+ & 9 & Saída apenas (ADC inoperante nesta placa) \\
\bottomrule
\end{tabular}
\end{center}

\subsubsection{Motores de passo}

\begin{center}
\small
\renewcommand{\arraystretch}{1.2}
\begin{tabular}{>{\bfseries}c c c c}
\toprule
\rowcolor{primary}
\textcolor{white}{Motor} & \textcolor{white}{Posição} & \textcolor{white}{STEP} & \textcolor{white}{DIR} \\
\midrule
A & topo & 15 & 16 \\
\rowstyle B & inferior-direito & 17 & 18 \\
C & inferior-esquerdo & 40 & 41 \\
\midrule
\multicolumn{4}{c}{\footnotesize ENABLE comum: \textbf{GPIO 10} (ativo em nível baixo)} \\
\bottomrule
\end{tabular}
\end{center}

\begin{infobox}[Escolha dos GPIOs]
Os pinos dos motores evitam deliberadamente os pinos da tela (4,5,8,9), os pinos
de \emph{strapping} (0,3,45,46) e os de USB (19,20). Há \emph{flags} de inversão
de sentido (\texttt{MOTOR\_A/B/C\_INVERT}) caso algum motor gire ao contrário no
\emph{homing}, sem precisar reconectar fios.
\end{infobox}

\subsection{Orientação física e convenção de eixos}

\begin{center}
\begin{tikzpicture}[font=\small]
  \draw[thick,fill=paperbg,rounded corners] (0,0) rectangle (5,3.4);
  \node at (2.5,3.75) {\scriptsize cabo FPC à direita $\rightarrow$};
  \node[primary] at (2.5,2.6) {\textbf{Motor A} (topo)};
  \node[primary] at (1.1,0.5) {\textbf{C}};
  \node[primary] at (3.9,0.5) {\textbf{B}};
  \node at (2.5,1.6) {\footnotesize MESA RESISTIVA};
  \draw[-{Stealth},accent,thick] (0.3,-0.4) -- (1.6,-0.4) node[right,scriptsize]{X (direita)};
  \draw[-{Stealth},secondary,thick] (-0.4,3.1) -- (-0.4,1.8) node[below,scriptsize,rotate=90,anchor=north]{};
  \node[scriptsize,secondary] at (-0.85,2.4) {Y};
\end{tikzpicture}
\end{center}

\begin{itemize}
  \item $x_{mm}=0$ na borda esquerda, $x_{mm}=187$ na direita (lado do cabo);
  \item $y_{mm}=0$ no topo, $y_{mm}=141$ na base;
  \item \emph{Setpoint} de equilíbrio: centro $(93{,}5,\ 70{,}5)$\,mm.
\end{itemize}

\subsection{Geometria 3RPS (parâmetros do firmware)}

\begin{center}
\small
\renewcommand{\arraystretch}{1.2}
\begin{tabular}{>{\ttfamily\bfseries}l c >{\RaggedRight\arraybackslash}p{7.5cm}}
\toprule
\rowcolor{primary}
\textcolor{white}{Constante} & \textcolor{white}{Valor} & \textcolor{white}{\normalfont\bfseries Significado} \\
\midrule
GEO\_D & 50,8\,mm & Raio da base (pivôs inferiores) \\
\rowstyle GEO\_E & 79,4\,mm & Raio da plataforma (juntas esféricas) \\
GEO\_F & 44,45\,mm & Comprimento do braço do motor \\
\rowstyle GEO\_G & 93,2\,mm & Comprimento da biela \\
GEO\_HZ & 108,0\,mm & Altura neutra da plataforma \\
\bottomrule
\end{tabular}
\end{center}

\begin{problembox}[Atenção: valores de referência]
Estes valores são \textbf{placeholders} herdados do projeto do Aaed Musa. Estão
marcados no código com a nota ``\emph{MEDIR no rig real e substituir}''. A
calibração geométrica fina é uma das pendências (Seção~9).
\end{problembox}

%==============================================================================
%  5. FIRMWARE E FERRAMENTAS
%==============================================================================
\section{Firmware, Ferramentas e Protocolo}

\subsection{Estrutura de módulos}

O firmware é escrito em C sobre o \textbf{ESP-IDF}. Cada módulo tem
responsabilidade única:

\begin{center}
\small
\renewcommand{\arraystretch}{1.25}
\begin{tabularx}{\textwidth}{>{\ttfamily\bfseries\color{primary}}l >{\RaggedRight\arraybackslash}X}
\toprule
main.c & Ponto de entrada; seleciona modo CONTROL ou MAPPING. \\
\rowstyle control.c & Laço de 100\,Hz, \emph{homing}, parser de comandos seriais, viés de nível (TRIM) e telemetria. \\
pid.c & PID por eixo: derivativo filtrado + anti-\emph{windup}. \\
\rowstyle kinematics.c & Cinemática inversa 3RPS (\texttt{machine\_theta}). \\
steppers.c & Gerador de passos STEP/DIR com perfil trapezoidal (timer 10\,kHz). \\
\rowstyle touch\_screen.c & Leitura resistiva 4 fios, média e conversão para mm. \\
cal\_store.c & Persistência em NVS (ganhos PID + calibração da tela). \\
\bottomrule
\end{tabularx}
\end{center}

\subsection{O laço de controle (100\,Hz)}

A cada ciclo, \texttt{control.c}:

\begin{enumerate}
  \item lê o toque (\texttt{touch\_read}, com detecção digital de presença e
  \textbf{rejeição de \emph{outliers}}) e mede o $\Delta t$ \emph{real} decorrido;
  \item com \textbf{presença \emph{debounced}}: se há bola, roda
  \texttt{pid\_update} para X e Y; sem bola, fica em IDLE (mesa nivelada) e reseta
  os PIDs;
  \item aplica os sinais de controle $(s_x,s_y)$ e soma o \textbf{viés de nível}
  $(\textsc{trim}_x,\textsc{trim}_y)$ — \emph{feedforward} para a base torta;
  \item converte a inclinação em ângulos via \texttt{apply\_tilt} $\rightarrow$
  \texttt{machine\_theta} e comanda os três motores (\texttt{steppers\_move\_to\_all});
  \item se habilitado, emite \texttt{POS}/\texttt{NOTOUCH}, \texttt{CTRL} e o
  \texttt{ERROR} (distância do centro, 1\,Hz).
\end{enumerate}

\begin{infobox}[\emph{Homing} em duas etapas (já implementado)]
Antes de equilibrar, a plataforma (1) \textbf{sobe} 15\,mm acima da altura de
trabalho — garantindo que o primeiro movimento dos braços seja sempre para cima,
longe da base — e (2) \textbf{desce} até nivelar ($n_x=n_y=0$). Resta apenas
\emph{validar} que a posição nivelada corresponde de fato à mesa plana
(pendência da Seção~9).
\end{infobox}

\subsection{Sensor resistivo}

A leitura faz a \textbf{média de 8 amostras de ADC por eixo}
(\param{TOUCH\_SAMPLES = 8}) e suaviza o resultado em mm com um filtro IIR
(\param{$\alpha$ = 0,45}). A \textbf{presença} da bola é detectada de forma
\textbf{digital} (pull-up, imune ao acoplamento capacitivo entre as camadas) e
\emph{independente da posição} — o que corrige a falha de não detectar nos cantos.
Uma \textbf{rejeição de \emph{outliers}} descarta saltos $>30$\,mm entre amostras
(glitch de leitura), evitando ``chutes'' na malha. A calibração mapeia o ADC bruto
para mm e o assistente \texttt{CAL} \textbf{detecta sozinho} \emph{swap}/\emph{flip}
a partir dos quatro cantos, \textbf{derivando o centro} deles (sem medi-lo com a bola).

\subsection{Ferramentas de PC}

\begin{center}
\small
\renewcommand{\arraystretch}{1.3}
\begin{tabularx}{\textwidth}{>{\ttfamily\bfseries\color{secondary}}l >{\RaggedRight\arraybackslash}X}
\toprule
bb\_tool & \textbf{Build / Flash / Monitor.} Menu interativo (raiz do projeto) que
carrega o ambiente ESP-IDF e oferece \emph{build}, \emph{fullclean}+\emph{set-target}
\texttt{esp32s3}, \emph{flash} e \emph{monitor} (porta padrão COM8). \\
\rowstyle calibrate.py & \textbf{Operação visual.} Calibração guiada dos 4 cantos,
bola ao vivo, ajuste de orientação e \emph{tuning} de PID em tempo real — persistido
no NVS sem recompilar. (A calibração e o PID também são feitos direto pelo terminal.) \\
PIDSimba.py & \textbf{Simulação física visual.} Bola virtual rolando sob
gravidade, três motores simulados com perfil trapezoidal; demonstra o laço PID
inteiro \emph{sem hardware}. Também espelha os motores reais (modo REAL). \\
\rowstyle sim\_pid.py & \textbf{Validador \emph{headless}} (só biblioteca padrão).
Roda a resposta em malha fechada contra o modelo da planta e imprime \emph{overshoot},
tempo de acomodação, erro residual e gráfico ASCII. \\
\bottomrule
\end{tabularx}
\end{center}

\subsection{Protocolo serial (USB-Serial-JTAG, 115200 baud)}

O console é o \textbf{USB-Serial-JTAG} (conector USB nativo do ESP32-S3), não a
UART externa. A referência completa de comandos está em \texttt{SERIAL\_COMMANDS.md}.

\subsubsection{Saída do firmware (telemetria com \texttt{SHOW}/\texttt{HIDDEN})}

\begin{center}
\small
\renewcommand{\arraystretch}{1.2}
\begin{tabularx}{\textwidth}{>{\ttfamily\footnotesize}l >{\RaggedRight\arraybackslash}X}
\toprule
\rowcolor{primary}
\textcolor{white}{\normalfont\bfseries Mensagem} & \textcolor{white}{\normalfont\bfseries Significado} \\
\midrule
POS,xr,yr,xmm,ymm & Bola detectada (ADC bruto + posição em mm). \\
\rowstyle NOTOUCH,xr,yr & Sem contato com a tela. \\
CTRL,x,y,sx,sy,nx,ny,thA,thB,thC & Estado do laço (pos., \emph{setpoint}, inclinação, ângulos). \\
\rowstyle ERROR,ex,ey,dist & Erro (distância do centro) a 1\,Hz, com \texttt{ERROR} ligado. \\
STATE,$\langle$PRONTO$\vert$ACTIVE$\vert$IDLE$\rangle$ & Transições de estado da bola. \\
\rowstyle GAINS,kp,ki,kd \quad SIGNS,sx,sy \quad TRIM,nx,ny & Ecos de configuração. \\
CAL,xmin,xmax,ymin,ymax,fx,fy,sw & Calibração da tela vigente. \\
\rowstyle SAVED,$\langle$pid$\vert$touch$\vert$trim$\rangle$ & Confirmação de gravação no NVS. \\
\bottomrule
\end{tabularx}
\end{center}

\subsubsection{Comandos de entrada}

\begin{center}
\small
\renewcommand{\arraystretch}{1.2}
\begin{tabularx}{\textwidth}{>{\ttfamily\footnotesize}l >{\RaggedRight\arraybackslash}X}
\toprule
\rowcolor{secondary}
\textcolor{white}{\normalfont\bfseries Comando} & \textcolor{white}{\normalfont\bfseries Efeito} \\
\midrule
HELP \quad ? & Lista de comandos \quad estado atual (GAINS/SIGNS/CAL). \\
\rowstyle SHOW / HIDDEN & Liga / desliga a telemetria. \\
ERROR & Imprime o erro (distância do centro) a cada 1\,s. \\
\rowstyle PID kp ki kd & Ajusta os três ganhos (\texttt{-} mantém o atual); \texttt{PID SAVE} grava. \\
KP/KI/KD v \quad SX/SY v & Um ganho por vez \quad sinal do controle por eixo. \\
\rowstyle TRIM \quad TRIM SAVE & Aprende / grava o viés de nível (base torta). \\
CAL & Assistente de calibração da tela (4 cantos; centro derivado). \\
\rowstyle STEPPER & Calibra o curso dos motores (MIN/MAX/MEIO). \\
PIDAUTO kp ki kd & Ajuste fino dos ganhos + auto-TRIM (parte dos ganhos atuais). \\
\rowstyle ELIPSE \quad DANCE & Varredura elíptica (sem bola) \quad dancinha de exibição. \\
CIRC & Gira a bola num círculo suave no centro (com bola; \texttt{PARAR} sai). \\
\rowstyle ZERO / ZEROCLR & Marca / limpa pontos mortos da tela. \\
\bottomrule
\end{tabularx}
\end{center}

\subsection{Modos de firmware}

\begin{center}
\small
\begin{tabularx}{\textwidth}{>{\ttfamily\bfseries}l >{\RaggedRight\arraybackslash}X}
\toprule
BB\_MODE\_CONTROL & (padrão) Laço PID 100\,Hz + cinemática + motores. \\
\rowstyle BB\_MODE\_MAPPING & Só lê a tela e emite telemetria — útil para depurar o sensor. \\
\bottomrule
\end{tabularx}
\end{center}

%==============================================================================
%  6. PROBLEMAS
%==============================================================================
\section{Problemas Enfrentados e Soluções}

Esta seção consolida as dificuldades reais registradas ao longo da construção
(\texttt{PROBLEMASBALL.md}), organizadas por natureza.

\subsection{Sensor: pinagem e leitura incorreta}

\begin{problembox}[Leitura errada da tela resistiva]
A pinagem da tela gerava \textbf{leitura de dados incorretos}, ou parte da mesa
simplesmente não respondia. O mapeamento elétrico dos quatro fios para os
GPIOs/canais de ADC não era óbvio, e a placa tinha uma particularidade: o pino
X+ (GPIO 9) \textbf{não funciona como ADC}, servindo apenas como saída.
\end{problembox}

\begin{solutionbox}
\begin{itemize}
  \item Mapeamento elétrico definitivo dos quatro fios (Seção~4.2), com X+ usado
  apenas como saída digital;
  \item leitura por \textbf{média de 8 amostras} por eixo mais filtro IIR
  ($\alpha=0{,}25$) para estabilizar o valor em mm;
  \item calibração in-situ dos quatro cantos, com \textbf{detecção automática}
  de \emph{swap}/\emph{flip} (\texttt{CAL APPLY}), persistida no NVS.
\end{itemize}
\end{solutionbox}

\subsection{Hardware de controle: a PCB}

\begin{problembox}[Necessidade de condensar a fiação]
Com a montagem em fios soltos espalhados, as ligações eram frágeis e propensas a
erro. Era necessário \textbf{condensar o hardware de controle em uma PCB} para
montar de forma adequada e confiável.
\end{problembox}

\begin{problembox}[Falha na fabricação da PCB]
A PCB foi enviada para fabricação, mas voltou com um \textbf{defeito: a placa
inteira em curto}. Não foi possível utilizá-la.
\end{problembox}

\begin{solutionbox}
A integração foi refeita em \textbf{placa perfurada}, com as trilhas formadas por
\textbf{soldagem manual} dos componentes. Mais trabalhoso, porém funcional —
permitiu validar o circuito e seguir com o projeto. A pinagem documentada
(Seção~4.2) fica como base para uma futura PCB corrigida.
\end{solutionbox}

\subsection{Estrutura mecânica}

\begin{problembox}[Impressão 3D sensível ao peso dos motores]
A estrutura impressa em 3D mostrou-se \textbf{relativamente sensível ao peso dos
três motores}, exigindo \textbf{reimpressão} de peças.
\end{problembox}

\begin{problembox}[Acesso a parafusos olhais]
Os \textbf{parafusos olhais} necessários às juntas eram de \textbf{difícil
acesso} no mercado.
\end{problembox}

\begin{solutionbox}
Reimpressão das peças críticas com maior robustez (mais \emph{infill}/reforço).
Quanto aos parafusos olhais, o que viabilizou a montagem foi o \textbf{professor
possuir} a peça — registrada agora na lista de materiais para aquisições futuras.
\end{solutionbox}

\subsection{Comportamento de controle: estabilidade, base torta e leituras espúrias}

\begin{problembox}[A mesa equilibrava e ``surtava do nada'']
Com a bola estável, a mesa eventualmente disparava para uma inclinação grande e
depois voltava. Duas causas: (i) o $K_d$ em uso estava muito abaixo do necessário,
deixando a malha \textbf{sub-amortecida} — e a planta é um \textbf{duplo
integrador} (Seção~3), em que o derivativo é \emph{quem} fornece o amortecimento;
(ii) uma \textbf{leitura espúria} isolada do painel virava um erro falso enorme que,
derivado, ``chutava'' a saída.
\end{problembox}

\begin{solutionbox}
Elevar o $K_d$ (amortecimento) e adicionar uma \textbf{rejeição de \emph{outliers}}
na leitura: saltos $>30$\,mm entre amostras (10\,ms) são descartados, a menos que
persistam (movimento real). A malha deixou de disparar com glitches.
\end{solutionbox}

\begin{problembox}[Bola estabiliza fora do centro e só converge devagar]
A estrutura inteira está levemente \textbf{torta}. Isso é uma \textbf{perturbação
constante (DC)}: a bola escorrega para um lado e só o termo \textbf{integral} zera
o erro de regime. Com $K_i$ pequeno, porém, o integral leva dezenas de segundos a
montar a correção — e \textbf{zera a cada bola} —, então o desvio reaparece sempre.
\end{problembox}

\begin{solutionbox}
Comando \textbf{\texttt{TRIM}}: a inclinação constante que o integral descobre é
transferida para um \textbf{\emph{feedforward} de nível} fixo, salvo no NVS e
aplicado já no boot (inclusive em repouso). A plataforma passa a \textbf{nascer
nivelada} em relação à gravidade; o integral cuida só do resíduo. É separar a
rejeição da perturbação DC (\emph{feedforward}) da regulação em torno dela (PID).
\end{solutionbox}

\begin{problembox}[Falsos toques e cantos não detectados]
O acoplamento capacitivo entre as camadas dava ``bola fantasma''; e exigir um
limiar de ADC para a presença \textbf{barrava os cantos}, onde a leitura de posição
cai legitimamente a $\approx 0$.
\end{problembox}

\begin{solutionbox}
Detecção de presença \textbf{digital} por pull-up (imune ao acoplamento) e
\textbf{independente da posição}; o ADC serve só para a coordenada. Calibração
\texttt{CAL} passou a usar \textbf{só os 4 cantos}, \textbf{derivando o centro}
deles (sem medi-lo com a bola, o que enviesava). Pontos presos residuais podem ser
mascarados com \texttt{ZERO}.
\end{solutionbox}

\subsection{Quadro-resumo}

\begin{center}
\small
\renewcommand{\arraystretch}{1.25}
\begin{tabularx}{\textwidth}{>{\bfseries}p{5.2cm} >{\RaggedRight\arraybackslash}X c}
\toprule
\rowcolor{primary}
\textcolor{white}{Problema} & \textcolor{white}{Solução} & \textcolor{white}{Criticidade} \\
\midrule
Leitura errada da tela & Pinagem definitiva + média/IIR + calibração & \textcolor{danger}{\textbf{Crítica}} \\
\rowstyle PCB em curto na fabricação & Soldagem em placa perfurada & \textbf{Alta} \\
Fiação solta / mau contato & Condensação na placa perfurada & \textbf{Alta} \\
\rowstyle Malha instável (``surta do nada'') & Mais $K_d$ (amortecimento) + rejeição de \emph{outliers} & \textbf{Alta} \\
Bola descentrada (base torta) & \emph{Feedforward} de nível (\texttt{TRIM}) & \textbf{Alta} \\
\rowstyle Falsos toques / cantos & Detecção digital; \texttt{CAL} de 4 cantos + \texttt{ZERO} & \textbf{Alta} \\
Impressão 3D frágil & Reimpressão reforçada & Média \\
\rowstyle Parafusos olhais difíceis & Fornecidos pelo professor & Média \\
\bottomrule
\end{tabularx}
\end{center}

%==============================================================================
%  7. CUSTOS
%==============================================================================
\section{Materiais e Custos}

\subsection{Recursos fornecidos (sem custo para o grupo)}

\begin{infobox}[Apoio institucional]
\begin{itemize}
  \item \textbf{Universidade (UNIFOR):} os três motores de passo.
  \item \textbf{Professor orientador:} a tela resistiva, parte dos parafusos
  (incluindo os olhais) e materiais de uso e consumo.
\end{itemize}
Esse apoio cobriu os itens de maior valor, viabilizando o projeto.
\end{infobox}

\subsection{Gastos diretos do grupo}

A tabela abaixo é o gasto \textbf{real} dos alunos até o momento.

\begin{center}
\renewcommand{\arraystretch}{1.3}
\begin{tabular}{>{\RaggedRight\arraybackslash}p{8cm} r}
\toprule
\rowcolor{primary}
\textcolor{white}{\bfseries Item} & \textcolor{white}{\bfseries Valor (R\$)} \\
\midrule
Filamento PLA (800\,g) & 90,00 \\
\rowstyle Drivers de motor — 3 unidades & 42,00 \\
Regulador de tensão & 17,00 \\
\rowstyle Porcas de inserção & 8,00 \\
Parafusos diversos & 10,00 \\
\rowstyle Placa perfurada & 15,90 \\
\midrule
\rowcolor{secondary}
\textcolor{white}{\bfseries TOTAL} & \textcolor{white}{\bfseries 182,90} \\
\midrule
\rowcolor{accent}
\textcolor{white}{\bfseries Por pessoa (8 integrantes)} & \textcolor{white}{\bfseries 22,86} \\
\bottomrule
\end{tabular}
\end{center}

\begin{infobox}[Leitura do custo]
O desembolso por aluno foi de apenas \textbf{R\$ 22,86}. O grupo arcou somente
com consumíveis e complementos (filamento, drivers, regulador, placa perfurada,
fixadores); os componentes de maior valor — motores e tela — vieram do apoio
institucional. O resultado é um projeto de controle completo a custo marginal
muito baixo por integrante.
\end{infobox}

%==============================================================================
%  8. VALIDAÇÃO
%==============================================================================
\section{Validação e Resultados}

\subsection{Validação teórica (simulação)}

\begin{itemize}
  \item \textbf{\texttt{sim\_pid.py} (modelo linear):} com os ganhos de alocação
  de polos, a resposta confirma \emph{overshoot} de $\approx 25$–$30\%$,
  $t_s\approx 3$\,s e erro estacionário nulo — coerente com $\zeta=0{,}636$ e
  $\omega_n=2{,}19$\,rad/s previstos analiticamente.
  \item \textbf{\texttt{PIDSimba.py} (física completa):} com os ganhos
  ``melhorados'', a bola converge ao centro \textbf{a partir dos quatro cantos},
  \emph{overshoot} $\approx 15\%$ e erro final $\approx 1$\,mm. \emph{A teoria
  fecha com o código.}
\end{itemize}

\subsection{Verificações de hardware}

\begin{center}
\small
\renewcommand{\arraystretch}{1.3}
\begin{tabularx}{\textwidth}{>{\bfseries}p{4.5cm} >{\RaggedRight\arraybackslash}X c}
\toprule
\rowcolor{primary}
\textcolor{white}{Item} & \textcolor{white}{Observação} & \textcolor{white}{Status} \\
\midrule
Leitura do sensor & Detecção digital + média/IIR + rejeição de \emph{outliers}; estável nos cantos & \textcolor{success}{\textbf{OK}} \\
\rowstyle Cinemática 3RPS & Inclinação mapeia para 3 ângulos; nivela no \emph{homing} & \textcolor{success}{\textbf{OK}} \\
Motores de passo & Perfil trapezoidal a 10\,kHz, DRV8825 em 1/16, sem perda de passos & \textcolor{success}{\textbf{OK}} \\
\rowstyle Laço 100\,Hz & Período e telemetria corretos & \textcolor{success}{\textbf{OK}} \\
Sinais $S_X/S_Y$ e equilíbrio & Realimentação negativa ($-1/-1$); bola converge ao centro & \textcolor{success}{\textbf{OK}} \\
\rowstyle Base torta (TRIM) & \emph{Feedforward} de nível compensa a estrutura & \textcolor{success}{\textbf{OK}} \\
Tuning fino do PID & Refinamento contínuo via \texttt{PID}/\texttt{PIDAUTO} & \textcolor{accent}{\textbf{Contínuo}} \\
\bottomrule
\end{tabularx}
\end{center}

\begin{infobox}[Como o tuning é feito no rig]
Direto pelo \textbf{terminal serial} (sem recompilar), com \texttt{ERROR} ligado
para ver a distância ao centro: (1) zere $K_i$ e $K_d$ e suba $K_p$ até a bola
\textbf{oscilar}; (2) adicione $K_d$ para \textbf{amortecer} (alvo $\zeta\approx 1$)
— \emph{passo crítico no duplo integrador}; (3) um \textbf{toque} de $K_i$ para o
desvio residual; (4) \texttt{TRIM} + \texttt{TRIM SAVE} para a base torta;
(5) \texttt{PID SAVE} grava no NVS. Ziegler–Nichols dá o chute inicial
($K_p\approx 0{,}6K_u$, $T_i\approx 0{,}5T_u$, $T_d\approx 0{,}125T_u$).
\end{infobox}

\begin{keybox}
\textbf{Configuração atual em operação:}\quad
$K_p = 6{,}0\times10^{-4}\;\;\; K_i = 3{,}0\times10^{-5}\;\;\; K_d = 1{,}0\times10^{-4}$
\quad·\quad sinais $-1/-1$ \quad·\quad CAL = \texttt{513,3297,355,3719,1,1,1}
\end{keybox}

%==============================================================================
%  9. PENDÊNCIAS
%==============================================================================
\section{Pendências e Trabalhos Futuros}

\subsection{Pendências técnicas imediatas}

Já \textbf{resolvidos} em operação: sinais $S_X/S_Y$ ($-1/-1$), microstepping em
1/16, calibração da tela (4 cantos) e do curso (\texttt{STEPPER}), detecção robusta
e compensação da base torta (\texttt{TRIM}). \emph{Restam:}

\begin{enumerate}
  \item \textbf{Medir a geometria real} (\texttt{GEO\_D/E/F/G/HZ}) com paquímetro
  e substituir os \emph{placeholders} em \texttt{control.c} — hoje o repouso é
  ancorado pelo curso calibrado (\texttt{STEPPER}), contornando os placeholders.
  \item \textbf{Refino contínuo de $K_p/K_i/K_d$} pelo terminal (\texttt{PID}) ou
  pelo auto-tune (\texttt{PIDAUTO}), persistindo no NVS.
  \item \textbf{Caracterizar a resposta} (overshoot/tempo de acomodação) com o
  monitor \texttt{ERROR} para comparar com a previsão analítica.
\end{enumerate}

\subsection{Melhorias futuras}

\begin{itemize}
  \item \textbf{PCB definitiva} (corrigindo o defeito de fabricação), a partir da
  pinagem já documentada.
  \item \textbf{Registro de telemetria} (cartão microSD ou \emph{log} no PC) para
  análise quantitativa da resposta.
  \item \textbf{Roteiro guiado} de montagem/calibração para as próximas turmas,
  reduzindo a curva de aprendizado.
\end{itemize}

%==============================================================================
%  10. CONCLUSÃO
%==============================================================================
\section{Conclusão}

O projeto entregou uma mesa equilibrista funcional, com controle PID digital
duplo embarcado no ESP32-S3, cinemática inversa 3RPS para três motores de passo e
sensor de posição resistivo calibrável em tempo de execução. A fundamentação
teórica — modelo de duplo integrador, projeto de ganhos por alocação de polos e
discretização com anti-\emph{windup} — está amarrada ao código real e validada em
simulação: a bola converge ao centro dos quatro cantos com \emph{overshoot} de
$\approx 15\%$.

A adaptação da abordagem de \textbf{dois servos} do TCC de Cordeiro para os
\textbf{três motores} deste rig mostrou-se consistente: a malha externa de
posição é preservada e a cinemática 3RPS realiza a inclinação. Os principais
obstáculos foram \textbf{práticos} — pinagem da tela, falha na fabricação da PCB
(contornada com placa perfurada), fragilidade da estrutura impressa e acesso a
componentes mecânicos — e foram superados com soluções documentadas.

O sistema está \textbf{operacional}; as pendências remanescentes (medição
geométrica, validação do \emph{homing} e \emph{tuning} fino) são refinamentos que
não comprometem o funcionamento básico e estão claramente mapeadas para a
continuidade. Com custo direto de apenas \textbf{R\$ 22,86 por integrante}, o
resultado é um banco de ensino completo e reaproveitável para futuras práticas de
controle.

%==============================================================================
%  REFERÊNCIAS
%==============================================================================
\section*{Referências}
\addcontentsline{toc}{section}{Referências}

\begin{enumerate}[label={[\arabic*]}]
  \item CORDEIRO, Giovanna Ferreira. \textit{Controle Clássico da Planta Ball
  Balancer.} Trabalho de Graduação, Engenharia Elétrica, UNESP/FEIS, Ilha
  Solteira, 2022.
  \item MUSA, Aaed. \textit{Ball Balancer.} Instructables. Disponível em:
  \url{https://www.instructables.com/Ball-Balancer/}.
  \item OGATA, Katsuhiko. \textit{Engenharia de Controle Moderno.} Prentice Hall.
  (2ª ordem, Routh–Hurwitz, alocação de polos, discretização.)
  \item FRANKLIN, G.; POWELL, J.; WORKMAN, M. \textit{Digital Control of Dynamic
  Systems.} (Euler/Tustin, projeto no domínio $z$.)
  \item ESPRESSIF SYSTEMS. \textit{ESP32-S3 Technical Reference Manual.}
  Disponível em: \url{https://www.espressif.com}.
  \item Documentação interna do projeto: \texttt{README.md} (montagem/operação) e
  \texttt{TEORIA.md} (controle digital).
\end{enumerate}

%==============================================================================
%  ANEXO
%==============================================================================
\appendix
\section{Anexo: Tabela-resumo de parâmetros}

\begin{center}
\small
\renewcommand{\arraystretch}{1.25}
\begin{tabularx}{\textwidth}{>{\ttfamily\bfseries}l c >{\RaggedRight\arraybackslash}X}
\toprule
\rowcolor{primary}
\textcolor{white}{Símbolo} & \textcolor{white}{Valor} & \textcolor{white}{\normalfont\bfseries Origem} \\
\midrule
A = (5/7)g & 7007\,mm/s² & Newton–Euler (esfera maciça) \\
\rowstyle T & 0,01\,s (100\,Hz) & \texttt{LOOP\_HZ} \\
$\tau_d$ & 0,03\,s & \texttt{PID\_D\_TAU} \\
\rowstyle $\alpha$ (filtro D) & 0,25 & $T/(\tau_d+T)$ \\
SIM\_KP & $1{,}08\times10^{-3}$ & alocação de polos \\
\rowstyle SIM\_KI & $3{,}0\times10^{-4}$ & ``PID melhorado'' \\
SIM\_KD & $9{,}0\times10^{-4}$ & ``PID melhorado'' \\
\rowstyle KP (firmware default) & $8{,}0\times10^{-4}$ & conservador, tunado no rig \\
KI (firmware default) & $2{,}0\times10^{-5}$ & $T_i = 40$\,s \\
\rowstyle KD (firmware default) & $1{,}2\times10^{-2}$ & freio com saturação proposital \\
$K_p,K_i,K_d$ (em operação) & $6{,}0{\times}10^{-4},\,3{,}0{\times}10^{-5},\,1{,}0{\times}10^{-4}$ & experimental, no rig (Seção~8) \\
TILT\_LIMIT & $\pm0{,}25$\,rad & saturação ($\approx 14{,}3°$) \\
\rowstyle $\zeta$ (projeto) & 0,636 & $PO\le 7{,}5\%$ \\
$\omega_n$ (projeto) & 2,19\,rad/s & $t_s\le 2{,}5$\,s \\
\rowstyle TOUCH\_SAMPLES & 8 & média de ADC por eixo \\
TOUCH\_IIR\_ALPHA & 0,45 & filtro de posição em mm \\
\rowstyle Salto p/ rejeição & 30\,mm/amostra & rejeição de \emph{outliers} \\
MICROSTEPS & 16 & DRV8825 (M0/M1/M2 = L/L/H) \\
\rowstyle SPEED\_BALANCE & 1800 passos/s & velocidade em equilíbrio \\
ACCEL\_BALANCE & 16000 passos/s² & aceleração em equilíbrio \\
\bottomrule
\end{tabularx}
\end{center}

\vfill
\begin{center}
{\color{primary}\rule{0.5\textwidth}{1pt}}\\[4pt]
{\small\color{inkdark}Ball Balancer 3RPS · Relatório Técnico · \today}
\end{center}

\end{document}
